% GENERATED FILE. Edit canonical lessons, then run npm run generate:books.
% canonical-source-hash: fnv1a64:f5e02140f68a5c55
% canonical-lessons: ZH-C231-zhongyi, ZH-W231-yi-medicine, ZH-C231-xiyi, ZH-W231-xi-west, ZH-C231-guahao, ZH-C231-hushi, ZH-C231-kaiyao

\chapter{Seeing a Doctor}
\label{ch:zh-seeing-a-doctor}
\hlchaptermodality{231}

\begin{chapteropening}
\textbf{By the end of this chapter:} \emph{I can choose Chinese or Western medicine, register at the clinic, talk about the nurse and say what the doctor gave me.}

Complete the last lesson of chapter 231: I can choose Chinese or Western medicine, register at the clinic, talk about the nurse and say what the doctor gave me.
\end{chapteropening}

\section[\texorpdfstring{\zh{中医} (zhōngyī)}{中医 (zhōngyī)}]{\zh{中医} (zhōngyī) — traditional Chinese medicine; a doctor of it}

\label{lesson:ZH-C231-zhongyi}



\begin{quote}
Before the new one: say the Mandarin for a nose, then the Mandarin for a leg.
\end{quote}

\subsection*{You'll want to know: \zh{中医}}
\textbf{\zh{中医}} — \emph{zhōngyī} — \textquotedblleft{}traditional Chinese medicine; a doctor of it\textquotedblright{}.

\begin{quote}
\textbf{\zh{我想看中医}} \emph{wǒ xiǎng kàn zhōngyī} — I'd like to see a Chinese-medicine doctor.
\end{quote}

The new character, in pieces: \zh{医} — seven strokes: a bar on top, five strokes inside it (a short fall, a short bar, a long bar, a long fall and a dot), then a stroke down the left that turns right along the foot. The next lesson puts it in your hand.

\begin{cousinweb}
\textbf{\zh{中医}} (\emph{zhōngyī}) is built from \zh{中}, as in \zh{中国}, and \zh{医} 'medicine', both inherited from Middle Chinese, \emph{trjuwng} and \emph{'i} (Guangyun, in Baxter's transcription), which glosses the second 'to heal'.
\end{cousinweb}

\begin{grammarlens}[title={What you've built}]
One more word for giving reasons.
\end{grammarlens}

\subsection*{Your turn}
\begin{itemize}
\raggedright
  \item \emph{Say it:} \emph{zhōngyī}
  \item \emph{Say it:} \emph{zhōngyī}, once more
  \item \emph{Say it:} \emph{wǒ xiǎng kàn zhōngyī} — I'd like to see a Chinese-medicine doctor.
  \item \emph{Recall:} say the Mandarin for a nose, then the Mandarin for a leg, then say \emph{zhōngyī} again
\end{itemize}

\begin{tcolorbox}[breakable,colback=teal!4,colframe=teal!35!black,title={Before you move on}]
What does \zh{中医} mean? (\textquotedblleft{}Traditional Chinese medicine; a doctor of it\textquotedblright{}.) Say it once more.
\end{tcolorbox}
\section[\texorpdfstring{\zh{医} (yī)}{医 (yī)}]{\zh{医} — a box open on the right, with five strokes inside}

\label{lesson:ZH-W231-yi-medicine}



\begin{quote}
The word first, then the character inside it.

\begin{itemize}
\raggedright
  \item \emph{Recall:} say \emph{zhōngyī} — \textbf{R1}, one lesson back, before it decays
\end{itemize}

\begin{quote}
\zh{头} \zh{中}
\end{quote}

\emph{Read it:} each one aloud, with nothing said first — \emph{tóu}, \emph{zhōng}
\end{quote}

\subsection*{Script — a box open on the right, with five strokes inside}
\hlblockfigure{\detokenize{figures/ZH-W231-yi-medicine-filmstrip.pdf}}{How \zh{医} is written, stroke by stroke}

\begin{quote}
\zh{医}
\end{quote}

\textbf{Two pieces: a box open on the right, and what sits inside.} The box takes two strokes, one first and one last.

\textbf{First, the top bar of the box.}

\textbf{Inside, five strokes.} A short stroke falling to the left, then a short bar. Then three as at the foot of \zh{头}: a long bar, a long stroke falling to the left through it, and a dot falling to the right.

\textbf{Last, the box closes:} a vertical down the left that turns and runs right along the foot.

1 + 5 + 1 = seven.

\subsection*{Your turn}
\begin{itemize}
\raggedright
  \item \emph{Write it:} \zh{医}, all seven, with the model in front of you
  \item \emph{Cover:} write it again from memory, and count the pieces, not the strokes
\end{itemize}

\begin{tcolorbox}[breakable,colback=teal!4,colframe=teal!35!black,title={Before you move on}]
How many strokes in \zh{医}? (\textbf{Seven}.) Which stroke comes last? (\textbf{The one down the left and along the foot}.) How many strokes sit inside? (\textbf{Five}.)

Source: \href{https://raw.githubusercontent.com/chanind/hanzi-writer-data/68d10a4b21150cae5e1ebbd223eed289cf32d90c/data/\%E5\%8C\%BB.json}{Hanzi Writer Data \zh{医}, pinned stroke paths and medians}.
\end{tcolorbox}
\section[\texorpdfstring{\zh{西医} (xīyī)}{西医 (xīyī)}]{\zh{西医} (xīyī) — western medicine; a doctor of it}

\label{lesson:ZH-C231-xiyi}



\begin{quote}
Before the new one: say the Mandarin for a leg, then the Mandarin for Chinese medicine.
\end{quote}

\subsection*{You'll want to know: \zh{西医}}
\textbf{\zh{西医}} — \emph{xīyī} — \textquotedblleft{}western medicine; a doctor of it\textquotedblright{}.

\begin{quote}
\textbf{\zh{我身体不好头也疼想看西医}} \emph{wǒ shēntǐ bù hǎo tóu yě téng xiǎng kàn xīyī} — I'm not well and my head hurts too; I'd like to see a doctor of western medicine.
\end{quote}

The new character, in pieces: \zh{西} — six strokes: a bar on top, a box of two strokes, two strokes inside it that start above it, and a bar along the foot. The next lesson puts it in your hand.

\begin{cousinweb}
\textbf{\zh{西医}} (\emph{xīyī}) is built from \zh{西} 'west' and \zh{医} 'medicine', both inherited from Middle Chinese, \emph{sej} and \emph{'i} (Guangyun, in Baxter's transcription), which explains the first, after an older book, as a bird settling on its nest when the sun is in the west.
\end{cousinweb}

\begin{grammarlens}[title={What you've built}]
One more word for giving reasons.
\end{grammarlens}

\subsection*{Your turn}
\begin{itemize}
\raggedright
  \item \emph{Say it:} \emph{xīyī}
  \item \emph{Say it:} \emph{xīyī}, once more
  \item \emph{Say it:} \emph{wǒ shēntǐ bù hǎo tóu yě téng xiǎng kàn xīyī} — I'm not well and my head hurts too; I'd like to see a doctor of western medicine.
  \item \emph{Recall:} say the Mandarin for a leg, then the Mandarin for Chinese medicine, then say \emph{xīyī} again
\end{itemize}

\begin{tcolorbox}[breakable,colback=teal!4,colframe=teal!35!black,title={Before you move on}]
What does \zh{西医} mean? (\textquotedblleft{}Western medicine; a doctor of it\textquotedblright{}.) Say it once more.
\end{tcolorbox}
\section[\texorpdfstring{\zh{西} (xī)}{西 (xī)}]{\zh{西} — a bar, a box and two strokes inside}

\label{lesson:ZH-W231-xi-west}



\begin{quote}
The word first, then the character inside it.

\begin{itemize}
\raggedright
  \item \emph{Recall:} say \emph{xīyī} — \textbf{R1}, one lesson back, before it decays
\end{itemize}

\begin{quote}
\zh{要} \zh{票}
\end{quote}

\emph{Read it:} each one aloud, with nothing said first — \emph{yào}, \emph{piào}
\end{quote}

\subsection*{Script — a bar, a box and two strokes inside}
\hlblockfigure{\detokenize{figures/ZH-W231-xi-west-filmstrip.pdf}}{How \zh{西} is written, stroke by stroke}

\begin{quote}
\zh{西}
\end{quote}

\textbf{One piece, six strokes.}

\textbf{A bar across the top.}

\textbf{The box, two strokes:} a vertical down the whole left side, then a bar that runs right and turns down.

\textbf{Inside, two strokes,} each starting above the box: a stroke down that falls a little to the left, then a stroke down that turns right along the inside.

\textbf{Last, a bar closes the foot.}

1 + 2 + 2 + 1 = six.

\subsection*{Your turn}
\begin{itemize}
\raggedright
  \item \emph{Write it:} \zh{西}, all six, with the model in front of you
  \item \emph{Cover:} write it again from memory, and count the pieces, not the strokes
\end{itemize}

\begin{tcolorbox}[breakable,colback=teal!4,colframe=teal!35!black,title={Before you move on}]
How many strokes in \zh{西}? (\textbf{Six}.) Where do the two inner strokes start? (\textbf{Above the box}.) What comes last? (\textbf{The bar along the foot}.)

Source: \href{https://raw.githubusercontent.com/chanind/hanzi-writer-data/68d10a4b21150cae5e1ebbd223eed289cf32d90c/data/\%E8\%A5\%BF.json}{Hanzi Writer Data \zh{西}, pinned stroke paths and medians}.
\end{tcolorbox}
\section[guàhào]{guàhào — to register at a clinic; to register at a hospital}

\label{lesson:ZH-C231-guahao}



\begin{quote}
Before the new one: say the Mandarin for Chinese medicine, then the Mandarin for western medicine.
\end{quote}

\subsection*{You'll want to know: guàhào}
\emph{guàhào} — \textquotedblleft{}to register at a clinic; to register at a hospital\textquotedblright{}.

\begin{quote}
\emph{kànbìng yǐqián yào guàhào} — Before you see the doctor, you have to register.
\end{quote}

Its characters come later; for now it is a word for the ear.

\begin{cousinweb}
\emph{guàhào} is built from \emph{guà} 'to hang up' and \zh{号} 'a number', both inherited from Middle Chinese, \emph{kweaH} and \emph{hawH} (Guangyun, in Baxter's transcription).
\end{cousinweb}

\begin{grammarlens}[title={What you've built}]
One more word for giving reasons.
\end{grammarlens}

\subsection*{Your turn}
\begin{itemize}
\raggedright
  \item \emph{Say it:} \emph{guàhào}
  \item \emph{Say it:} \emph{guàhào}, once more
  \item \emph{Say it:} \emph{kànbìng yǐqián yào guàhào} — Before you see the doctor, you have to register.
  \item \emph{Recall:} say the Mandarin for Chinese medicine, then the Mandarin for western medicine, then say \emph{guàhào} again
\end{itemize}

\begin{tcolorbox}[breakable,colback=teal!4,colframe=teal!35!black,title={Before you move on}]
What does guàhào mean? (\textquotedblleft{}To register at a clinic; to register at a hospital\textquotedblright{}.) Say it once more.
\end{tcolorbox}
\section[hùshi]{hùshi — nurse}

\label{lesson:ZH-C231-hushi}



\begin{quote}
Before the new one: say the Mandarin for western medicine, then the Mandarin for to register at a clinic.
\end{quote}

\subsection*{You'll want to know: hùshi}
\emph{hùshi} — \textquotedblleft{}nurse\textquotedblright{}.

\begin{quote}
\emph{tā shì zhèr de hùshi} — She's a nurse here.
\end{quote}

Its characters come later; for now it is a word for the ear.

\begin{cousinweb}
\emph{hùshi} is built from \emph{hù} 'to protect' and \emph{shi} 'a scholar', both inherited from Middle Chinese, \emph{huH} and \emph{dzriX} (Guangyun, in Baxter's transcription), which glosses the first 'to save' and 'to help'; in the word the second loses its tone.
\end{cousinweb}

\begin{grammarlens}[title={What you've built}]
One more word for giving reasons.
\end{grammarlens}

\subsection*{Your turn}
\begin{itemize}
\raggedright
  \item \emph{Say it:} \emph{hùshi}
  \item \emph{Say it:} \emph{hùshi}, once more
  \item \emph{Say it:} \emph{tā shì zhèr de hùshi} — She's a nurse here.
  \item \emph{Recall:} say the Mandarin for western medicine, then the Mandarin for to register at a clinic, then say \emph{hùshi} again
\end{itemize}

\begin{tcolorbox}[breakable,colback=teal!4,colframe=teal!35!black,title={Before you move on}]
What does hùshi mean? (\textquotedblleft{}Nurse\textquotedblright{}.) Say it once more.
\end{tcolorbox}
\section[kāiyào]{kāiyào — to prescribe medicine}

\label{lesson:ZH-C231-kaiyao}



\begin{quote}
Before the new one: say the Mandarin for to register at a clinic, then the Mandarin for a nurse.
\end{quote}

\subsection*{You'll want to know: kāiyào}
\emph{kāiyào} — \textquotedblleft{}to prescribe medicine\textquotedblright{}.

\begin{quote}
\emph{zhōngyī yě huì kāiyào} — Chinese-medicine doctors prescribe medicine too.
\end{quote}

Its characters come later; for now it is a word for the ear.

\begin{cousinweb}
\emph{kāiyào} is built from \zh{开} 'to open' and \emph{yào} 'medicine', both inherited from Middle Chinese, \emph{khoj} and \emph{yak} (Guangyun, in Baxter's transcription); the sense 'to write out a prescription' is not given there.
\end{cousinweb}

\begin{grammarlens}[title={What you've built}]
One more word for giving reasons.
\end{grammarlens}

\subsection*{Your turn}
\begin{itemize}
\raggedright
  \item \emph{Say it:} \emph{kāiyào}
  \item \emph{Say it:} \emph{kāiyào}, once more
  \item \emph{Say it:} \emph{zhōngyī yě huì kāiyào} — Chinese-medicine doctors prescribe medicine too.
  \item \emph{Recall:} say the Mandarin for to register at a clinic, then the Mandarin for a nurse, then say \emph{kāiyào} again
\end{itemize}

\begin{tcolorbox}[breakable,colback=teal!4,colframe=teal!35!black,title={Before you move on}]
What does kāiyào mean? (\textquotedblleft{}To prescribe medicine\textquotedblright{}.) Say it once more.
\end{tcolorbox}
